% !TEX root = ../main.tex
\section{Producto propuesto}
El entregable principal de este proyecto no se limiuta a la simple configuracion de una herramienta de software, sino que
constituye un modelo arquitectonico de referencia para la modernizacion de la entrega de valor en sistemas legados. La
solucion transformara el actual paradigma de despliegue artesanal, fragmentado y vulnerable, en una cadena de suministro
de software predecible, auditable y badsada en principios de Infraestructura como Codigo (IaC).

Las innovaciones tecnologicas y operativas que caracterizan a esta propuesta incluyen:
\begin{itemize}
    \item \textbf{Estandar de Inmutabilidad:} Transacion hacia contenedores esatndarizados que erradican la divergencia de 
    configuracion entre los entornos de desarrollo local y los servidores en vivo.
    \item \textbf{Orquetacion Declarativa en Entornos Restringidos:} Implementacion de un flujo de trabajo CI/CD integral operado
    enteramente en \emph{on-premise}, demostrando la viabilidad de la cultura DevOps sin depender de proveedores de nube publica.
    \item \textbf{Entornos de Validacion de Alta Fidelidad:} Resolucion del deficit de pruebas automatizadas mediante la instacion
    de replicas aisladas del sistema. Estas instancias, poblados con datos estrcuturales reales, habilitaran la verificacion funcional
    humana en condiciones indeticas a produccion antes de la fusion en codigo con la rama principal.
    \item \textbf{Transaccionalidad en Esquemas Heterogeneos:} Automatizacion rigurosa de las migraciones de base de datos
    que sincroniza los cambios estrdcuturales entre motores (MySQL y MongoDB) como un unico evento atomico de despliegue.
\end{itemize}

\begin{figure}[htbp]
\centering
\framebox{\parbox{0.8\textwidth}{\centering
\vspace{3cm}
\textbf{Diagrama de Arquitectura de la Solución (Placeholder)} \
\small\textit{Aquí irá un diagrama de arquitectura mostrando la transición de código desde el repositorio local hasta la producción, destacando el rol del Runner autoalojado, la creación de imágenes y el aprovisionamiento de las bases duales en el entorno de prueba.}
\vspace{3cm}
}}
\caption{Arquitectura del modelo de entrega continua propuesto para el sistema heredado.}
\label{fig:arquitectura_solucion}
\end{figure}

\section{Metodología}

[Pasos o fases para alcanzar el objetivo, cada uno con definición y
descripción de cómo se realizará.]

\section{Programa de actividades}

[Cronograma de actividades por periodo.]

\section{Resultados esperados}

[Qué se espera obtener al terminar.]
